%% 第4巻 巻末（こたえ・記号カード・定数と単位・公式・チェック表・つぎの巻へ・さくいん・奥付）（第13回編集委員会）
%% 第3巻の巻末と同じ形。記号カードの一覧・書くときのちゅうい・定数と単位・公式・表計算・つぎの巻へを、第4巻の本文に合わせて書いた
%% 第4巻はルビなし（\R は使わない）。字の色は bkTheme でなく \bkTT
\makeatletter
%% ---------------- こたえ ----------------
\bkAnswerPages

%% ---------------- 記号カード ----------------
%% 前の巻の記号で、この巻でも使うもの（この巻の章の \bkkigocard は自動でうしろに集まる。θ・ω などは章のカードの文になる）
\bkKigoPages{%
\bkkigorow{$<$\ \ $>$\ \ $\leqq$\ \ $\geqq$}{不等号}{大きさをくらべる記号。$\leqq$ は「以下」、$\geqq$ は「以上」}{2-2}%
\bkkigorow{$\Delta$}{デルタ}{変化量。「あと$-$まえ」。三角形の印 $\triangle$（4-1）とはちがう}{2-3}%
\bkkigorow{$x^2$}{エックスの2乗}{$x\times x$。$x^3$ は $x\times x\times x$}{2-7}%
\bkkigorow{$\sqrt{\phantom{00}}$}{ルート}{2乗すると中の数になる正の数。$\sqrt{9}=3$}{2-8}%
\bkkigorow{$\pi$}{パイ}{円周率。円周÷直径。$3.14159\cdots$}{2-9}%
\bkkigorow{$x_1$}{エックス・いち}{1番めの $x$。右下の数字は名札で、かけ算ではない}{3-2}%
\bkkigorow{$\lvert x\rvert$}{絶対値エックス}{0からの距離。$\lvert -3\rvert=3$}{3-2}%
\bkkigorow{$f(x)$}{エフ・エックス}{関数 $f$ に $x$ を入れた値}{3-3}%
\bkkigorow{$\approx$}{ニアリーイコール}{ほぼ等しい。$\sqrt2\approx1.414$}{3-17}%
}

%% 記号カードのつづき：ギリシャ文字（この巻から、角の θ・ω・α・φ などがふえる）
\def\bkbkhead#1{\par\addvspace{4mm}\noindent{\sffamily\bfseries\large\color{\bkTT}#1}\par\nobreak\vspace{2mm}}
\par\noindent\begin{minipage}{\linewidth}
\bkbkhead{ギリシャ文字の読み方}
{\small この巻から、角を表す $\theta$・$\alpha$・$\varphi$ や、角速度の $\omega$ など、ギリシャ文字がふえます。大文字・小文字・読み方と、よく使われる意味です。}\par\vspace{2mm}
\begin{bkcenter}
\footnotesize
\begin{tblr}{colspec={Q[c,m,wd=7mm] Q[c,m,wd=7mm] Q[l,m,wd=17mm] Q[l,m,wd=30mm] Q[c,m,wd=7mm] Q[c,m,wd=7mm] Q[l,m,wd=17mm] Q[l,m,wd=30mm]}, row{1}={bg=bkVol!14, font=\sffamily\bfseries\footnotesize}, hlines={0.4pt,bkLine}, vline{5}={0.8pt,bkVol}, rowsep=0.8pt, colsep=2.5pt}
大 & 小 & 読み方 & よく使う意味 & 大 & 小 & 読み方 & よく使う意味 \\
$\mathrm{A}$ & $\alpha$ & アルファ & 角（4-9） & $\mathrm{N}$ & $\nu$ & ニュー & 振動数 \\
$\mathrm{B}$ & $\beta$ & ベータ & 角 & $\Xi$ & $\xi$ & クシー & （変数） \\
$\Gamma$ & $\gamma$ & ガンマ & （係数） & $\mathrm{O}$ & $o$ & オミクロン & ― \\
$\Delta$ & $\delta$ & デルタ & 変化量、小さな量 & $\Pi$ & $\pi$ & パイ & 円周率 \\
$\mathrm{E}$ & $\varepsilon$ & イプシロン & 小さな数 & $\mathrm{P}$ & $\rho$ & ロー & 密度 \\
$\mathrm{Z}$ & $\zeta$ & ゼータ & （係数） & $\Sigma$ & $\sigma$ & シグマ & 和、ばらつき \\
$\mathrm{H}$ & $\eta$ & イータ & 効率 & $\mathrm{T}$ & $\tau$ & タウ & トルク \\
$\Theta$ & $\theta$ & シータ & 角（4-2） & $\Upsilon$ & $\upsilon$ & ウプシロン & ― \\
$\mathrm{I}$ & $\iota$ & イオタ & ― & $\Phi$ & $\varphi$ & ファイ & 球座標の角（4-15） \\
$\mathrm{K}$ & $\kappa$ & カッパ & （係数） & $\mathrm{X}$ & $\chi$ & カイ & ― \\
$\Lambda$ & $\lambda$ & ラムダ & 波長 & $\Psi$ & $\psi$ & プサイ & （角） \\
$\mathrm{M}$ & $\mu$ & ミュー & 摩擦係数（4-6） & $\Omega$ & $\omega$ & オメガ & 角速度（4-8） \\
\end{tblr}
\end{bkcenter}
{\footnotesize 大文字が英語のアルファベットと同じ形の文字（A、B、E など）は、記号としてはほとんど使いません。$\varphi$ は $\phi$ と書くこともあります。}
\end{minipage}

%% 記号カードのつづき：手で書くときのちゅうい（各章の「まちがえやすいところ」から）
\clearpage
\bkbkhead{記号を書くときのちゅうい}
{\small 手で書くと、ほかの字や記号とまぎれやすいものと、読みまちがえやすい書き方です。}\par\vspace{2mm}
\begin{bkcenter}
\begin{tblr}{colspec={Q[c,m,wd=19mm] Q[l,m,wd=22mm] Q[l,m,wd=85mm]}, row{1}={bg=bkVol!14, font=\sffamily\bfseries\footnotesize}, hlines={0.4pt,bkLine}, rowsep=0.8pt, cells={font=\small}}
記号 & まぎれやすい & 書くとき・読むときのコツ \\
$\sin\theta$ & $\sin\times\theta$ & $\sin$ と $\theta$ で1つの数（角 $\theta$ のサイン）。かけ算ではない（4-6） \\
$\sin^{-1}0.5$ & $\dfrac{1}{\sin0.5}$ & 「$\sin$ が0.5になる角」で\ang{30}。$-1$ は逆数の意味ではない（4-6） \\
$\sin^2\theta$ & $\sin(\theta^2)$ & $(\sin\theta)^2$ のこと。2は $\sin$ のすぐ右上に書く（4-8） \\
DEG・RAD & ― & 電卓の角の単位を先に確かめる（度は DEG、ラジアンは RAD）。表計算はラジアン（4-2） \\
$\theta$ & 0、O & たて長の丸の、まん中に横線を1本（4-2） \\
$\omega$ & $w$ & 下を丸く2つ。$w$ のようにとがらせない（4-8） \\
$\angle\mathrm{ABC}$ & ― & まん中に書いた B が頂点の角（4-1） \\
$\triangle$ & $\Delta$ & 三角形の印。変化量の $\Delta$（2-3）とは別の記号（4-1） \\
$\sojisym$\ \ $\equiv$ & $\sim$、$=$ & $\sojisym$ は相似（S を横にした形。なみの ～ とはちがう）、横線3本の $\equiv$ は合同。頂点は対応する順に（4-3） \\
$\vec a$ & $a$ & 矢印がないと、ただの数（大きさ）。$\overrightarrow{\mathrm{AB}}$ は A から B へ（4-10） \\
$\symbf a$ & $a$ & 太字はベクトル、細い字はその大きさ。ノートでは $\vec a$ と書く（4-13） \\
$\vec a\cdot\vec b$ & $\vec a\times\vec b$ & 内積の点はまん中に。答えは数。外積とは別の計算（4-12） \\
$\vec a\times\vec b$ & $\vec b\times\vec a$ & 順番を入れかえると反対向き。$\vec b\times\vec a=-\vec a\times\vec b$（4-14） \\
$\odot$\ \ $\otimes$ & ― & 丸に点は手前へ出る向き、丸に $\times$ はおくへ入る向き（4-14） \\
\end{tblr}
\end{bkcenter}

%% ---------------- 定数・単位・公式 ----------------
\bkbacktitle{よく使う定数・単位・公式}
\bksubsection{1}{この巻で使う単位}
\begin{bkcenter}
\small
\begin{tblr}{colspec={Q[l,m,wd=24mm] Q[l,m,wd=34mm] Q[l,m,wd=49mm] Q[c,m,wd=14mm]}, row{1}={bg=bkVol!14, font=\sffamily\bfseries\footnotesize}, hlines={0.4pt,bkLine}, rowsep=1.4pt}
量 & 単位（読み方） & ほかの単位との関係 & 章 \\
長さ & m、cm、mm、km & $\qty{1}{m}=\qty{100}{cm}$、$\qty{1}{km}=\qty{1000}{m}$ & 4-3 \\
 & in（インチ） & $\qty{1}{in}=\qty{2.54}{cm}$ & 4-4 \\
 & nm（ナノメートル） & \qty{1}{nm}は10億分の1\,m & 4-13 \\
体積 & m$^3$、L、cm$^3$ & $\qty{1}{m^3}=\qty{1000}{L}$、$\qty{1}{L}=\qty{1000}{cm^3}$ & 4-5 \\
角 & °（度）、rad（ラジアン） & $\pi$\,rad＝\ang{180}、\ $1^{\circ}=\dfrac{\pi}{180}$\,rad、\ 1\,rad $\approx$ \ang{57.3} & 4-2 \\
角速度 & rad/s & 1秒に回る角。$\omega=\dfrac{\theta}{t}$ & 4-8 \\
振動数 & Hz（ヘルツ） & 1秒に何回。$f=\dfrac{1}{T}$、\ $\omega=2\pi f$ & {4-8\\4-9} \\
力 & N（ニュートン） & 重力 $mg$（$g\approx\qty{9.8}{N/kg}$） & 4-6 \\
仕事 & J（ジュール） & $\qty{1}{J}=1$\,N·m & 4-12 \\
力のモーメント & N·m & 力×うでの長さ（トルク） & 4-14 \\
磁場・電荷 & T（テスラ）、C（クーロン） & 磁場から受ける力 $q\vec v\times\vec B$ & 4-14 \\
\end{tblr}
\end{bkcenter}

\bksubsection{2}{よく使う数}
\begin{bkcenter}
\small
\begin{tblr}{colspec={Q[l,m,wd=42mm] Q[l,m,wd=60mm] Q[c,m,wd=20mm]}, row{1}={bg=bkVol!14, font=\sffamily\bfseries\footnotesize}, hlines={0.4pt,bkLine}, rowsep=1pt}
名前 & 値 & 章 \\
円周率 $\pi$ & 3.14159… & 2-9・4-2 \\
$\sqrt2$、\ $\sqrt3$ & 1.414…、\ 1.732… & 2-8・4-6 \\
1\,rad & 約\ang{57.3}（1周は約6.28\,rad） & 4-2 \\
重力の強さ $g$ & 約\qty{9.8}{N/kg}（\unit{m/s^2} と同じ単位） & 4-6 \\
1インチ & \qty{2.54}{cm} & 4-4 \\
\end{tblr}
\end{bkcenter}

\bksubsection{3}{図形の公式}
{\small
\begin{longtblr}{colspec={Q[l,m,wd=28mm] Q[l,m,wd=84mm] Q[c,m,wd=12mm]}, rowhead=1, row{1}={bg=bkVol!14, font=\sffamily\bfseries\footnotesize}, hlines={0.4pt,bkLine}, rowsep=1.6pt}
公式・きまり & 式 & 章 \\
平行線と角 & 対頂角は等しい。平行な2本の直線では、同位角は等しく、錯角は等しい & 4-1 \\
内角・外角 & 三角形の内角の和は\ang{180}。$n$ 角形の内角の和＝$\ang{180}\times(n-2)$。外角の和は\ang{360} & 4-1 \\
円周・弧度法 & 円周＝$2\pi r$。角 $\theta$（rad）＝弧の長さ÷半径。1周は $2\pi$\,rad（\ang{360}） & 4-2 \\
弧とおうぎ形 & 弧の長さ $s=r\theta$、\ おうぎ形の面積 $S=\dfrac{1}{2}rs=\dfrac{1}{2}r^2\theta$（$\theta$ はラジアン） & 4-2 \\
円周角・接線 & 円周角は中心角の半分。直径に対する円周角は\ang{90}。接線は接点で半径に垂直 & 4-2 \\
相似 & 相似比が $1:k$ なら、長さの比 $1:k$、面積の比 $1:k^2$、体積の比 $1:k^3$ & 4-3 \\
相似条件 & ① 3辺の比がすべて等しい\quad ② 2辺の比とその間の角が等しい\quad ③ 2つの角が等しい & 4-3 \\
平行線と比 & DE∥BC なら\quad \mbox{AD : AB＝AE : AC＝DE : BC} & 4-3 \\
三平方の定理 & 直角三角形で、斜辺を $c$ とすると\quad $a^2+b^2=c^2$（逆：$a^2+b^2=c^2$ なら $c$ の向かいの角は直角） & 4-4 \\
三角定規 & $1:1:\sqrt2$（\ang{45}）、\ $1:2:\sqrt3$（\ang{30}・\ang{60}） & 4-4 \\
2点の距離・円 & $d=\sqrt{(x_2-x_1)^2+(y_2-y_1)^2}$。原点が中心、半径 $r$ の円は $x^2+y^2=r^2$ & 4-4 \\
直方体の対角線 & $\sqrt{a^2+b^2+c^2}$ & 4-4 \\
面積 & 平行四辺形＝底辺×高さ、三角形＝底辺×高さ÷2、台形＝（上底＋下底）×高さ÷2、ひし形＝対角線×対角線÷2 & 4-5 \\
円・おうぎ形 & 円 $\pi r^2$、\ おうぎ形 $\dfrac{1}{2}\times$弧の長さ$\times r$ & 4-5 \\
体積・表面積 & 柱＝底面積×高さ、すい＝底面積×高さ×$\dfrac{1}{3}$、球の体積 $\dfrac{4}{3}\pi r^3$、球の表面積 $4\pi r^2$ & 4-5 \\
\end{longtblr}\par}

\bksubsection{4}{三角比と三角関数の公式}
{\small
\begin{longtblr}{colspec={Q[l,m,wd=28mm] Q[l,m,wd=84mm] Q[c,m,wd=12mm]}, rowhead=1, row{1}={bg=bkVol!14, font=\sffamily\bfseries\footnotesize}, hlines={0.4pt,bkLine}, rowsep=1.6pt}
公式・きまり & 式 & 章 \\
三角比の定義 & $\sin\theta=\dfrac{\text{対辺}}{\text{斜辺}}$\quad $\cos\theta=\dfrac{\text{隣辺}}{\text{斜辺}}$\quad $\tan\theta=\dfrac{\text{対辺}}{\text{隣辺}}$（坂の傾き＝$\tan\theta$） & 4-6 \\
相互関係 & $\tan\theta=\dfrac{\sin\theta}{\cos\theta}$、\ $\sin^2\theta+\cos^2\theta=1$、\ $\sin(\ang{90}-\theta)=\cos\theta$ & {4-6\\4-8} \\
鈍角 & $\sin(\ang{180}-\theta)=\sin\theta$、\ $\cos(\ang{180}-\theta)=-\cos\theta$（鈍角の $\cos$ は負） & 4-6 \\
斜面の重力 & 斜面にそって $mg\sin\theta$、斜面へ押しつける向きに $mg\cos\theta$ & 4-6 \\
三角形の面積 & $S=\dfrac{1}{2}ab\sin C$（2辺とそのあいだの角） & 4-7 \\
正弦定理 & $\dfrac{a}{\sin A}=\dfrac{b}{\sin B}=\dfrac{c}{\sin C}=2R$（$R$ は外接円の半径） & 4-7 \\
余弦定理 & $c^2=a^2+b^2-2ab\cos C$、\ $\cos C=\dfrac{a^2+b^2-c^2}{2ab}$ & 4-7 \\
ロボットのうで & $x=L_1\cos\theta_1+L_2\cos(\theta_1+\theta_2)$、\ $y=L_1\sin\theta_1+L_2\sin(\theta_1+\theta_2)$ & 4-7 \\
単位円・一般角 & 単位円の上の点は $(\cos\theta,\ \sin\theta)$。$\theta+\ang{360}\times n$ は $\theta$ と同じ位置（周期\ang{360}＝$2\pi$） & 4-8 \\
\ang{90}回す & $(x,\ y)\to(-y,\ x)$。\ $\cos(\theta+\ang{90})=-\sin\theta$、\ $\sin(\theta+\ang{90})=\cos\theta$ & 4-8 \\
回る点 & $\theta=\omega t$、\ $(r\cos\omega t,\ r\sin\omega t)$、\ $T=\dfrac{2\pi}{\omega}$、\ $f=\dfrac{1}{T}$、\ $v=r\omega$ & 4-8 \\
加法定理 & $\sin(a\pm b)=\sin a\cos b\pm\cos a\sin b$、\ $\cos(a\pm b)=\cos a\cos b\mp\sin a\sin b$ & 4-9 \\
2倍角 & $\sin2\theta=2\sin\theta\cos\theta$、\ $\cos2\theta=1-2\sin^2\theta=2\cos^2\theta-1$ & 4-9 \\
2乗を1次に & $\sin^2\theta=\dfrac{1-\cos2\theta}{2}$、\ $\cos^2\theta=\dfrac{1+\cos2\theta}{2}$ & 4-9 \\
合成 & $a\sin\theta+b\cos\theta=\sqrt{a^2+b^2}\,\sin(\theta+\alpha)$ & 4-9 \\
積を和に & $2\sin a\cos b=\sin(a+b)+\sin(a-b)$、\ $2\sin a\sin b=\cos(a-b)-\cos(a+b)$ & 4-9 \\
和を積に & $\sin A+\sin B=2\sin\dfrac{A+B}{2}\cos\dfrac{A-B}{2}$ & 4-9 \\
\end{longtblr}\par}

\vspace{2mm}
\noindent\begin{minipage}{\linewidth}
\noindent{\sffamily\bfseries\small\color{\bkTT}よく使う角の三角比}（4-6）\par\vspace{1mm}
\begin{bkcenter}
\begin{tblr}{colspec={Q[c,m,wd=12mm] *{9}{Q[c,m,wd=10.4mm]}}, row{1}={bg=bkVol!14, font=\sffamily\bfseries\footnotesize}, hlines={0.4pt,bkLine}, rowsep=1.6pt, colsep=2pt, cells={font=\small}}
$\theta$ & \ang{0} & \ang{30} & \ang{45} & \ang{60} & \ang{90} & \ang{120} & \ang{135} & \ang{150} & \ang{180} \\
$\sin\theta$ & 0 & $\frac{1}{2}$ & $\frac{\sqrt2}{2}$ & $\frac{\sqrt3}{2}$ & 1 & $\frac{\sqrt3}{2}$ & $\frac{\sqrt2}{2}$ & $\frac{1}{2}$ & 0 \\
$\cos\theta$ & 1 & $\frac{\sqrt3}{2}$ & $\frac{\sqrt2}{2}$ & $\frac{1}{2}$ & 0 & $-\frac{1}{2}$ & $-\frac{\sqrt2}{2}$ & $-\frac{\sqrt3}{2}$ & $-1$ \\
$\tan\theta$ & 0 & $\frac{\sqrt3}{3}$ & 1 & $\sqrt3$ & なし & $-\sqrt3$ & $-1$ & $-\frac{\sqrt3}{3}$ & 0 \\
\end{tblr}
\end{bkcenter}
{\footnotesize ラジアンでは、\ang{30}＝$\dfrac{\pi}{6}$、\ang{45}＝$\dfrac{\pi}{4}$、\ang{60}＝$\dfrac{\pi}{3}$、\ang{90}＝$\dfrac{\pi}{2}$、\ang{180}＝$\pi$（4-2）。$\sqrt2\approx1.414$、$\sqrt3\approx1.732$ なので、$\dfrac{\sqrt2}{2}\approx0.707$、$\dfrac{\sqrt3}{2}\approx0.866$。}
\end{minipage}

\bksubsection{5}{ベクトルと座標の公式}
{\small
\begin{longtblr}{colspec={Q[l,m,wd=28mm] Q[l,m,wd=84mm] Q[c,m,wd=12mm]}, rowhead=1, row{1}={bg=bkVol!14, font=\sffamily\bfseries\footnotesize}, hlines={0.4pt,bkLine}, rowsep=1.6pt}
公式・きまり & 式 & 章 \\
和・実数倍 & 矢印をつなぐ。$\overrightarrow{\mathrm{AB}}+\overrightarrow{\mathrm{BC}}=\overrightarrow{\mathrm{AC}}$。$k\vec a$ は長さ $|k|$ 倍（$k<0$ なら反対向き） & 4-10 \\
差 & $\vec a-\vec b=\vec a+(-\vec b)$。\ $\overrightarrow{\mathrm{OA}}-\overrightarrow{\mathrm{OB}}=\overrightarrow{\mathrm{BA}}$（$\vec b$ の先から $\vec a$ の先へ） & 4-10 \\
速度の変化 & $\Delta\vec v=\vec v_{\text{あと}}-\vec v_{\text{まえ}}$、\ 相対速度 $\vec v_{\mathrm{AB}}=\vec v_{\mathrm{B}}-\vec v_{\mathrm{A}}$ & 4-10 \\
成分 & $\vec a=(a_x,\ a_y)$、\ $\overrightarrow{\mathrm{AB}}=(x_2-x_1,\ y_2-y_1)$。和・差・実数倍は成分ごとに計算 & 4-11 \\
大きさ・単位ベクトル & $|\vec a|=\sqrt{a_x^{\,2}+a_y^{\,2}}$、\ $\vec e=\dfrac{\vec a}{|\vec a|}$ & 4-11 \\
力の分解・つり合い & $\vec F=(F\cos\theta,\ F\sin\theta)$。つり合いは $x$ 成分の和 $=0$、$y$ 成分の和 $=0$ & 4-11 \\
中点・重心 & 中点 $\dfrac{\vec a+\vec b}{2}$、\ 重心 $\dfrac{\vec a+\vec b+\vec c}{3}$、\ $\dfrac{m_1\vec r_1+m_2\vec r_2}{m_1+m_2}$ & 4-11 \\
仕事 & $W=Fs\cos\theta=\vec F\cdot\vec s$（力と動く向きが直角なら0） & 4-12 \\
内積 & $\vec a\cdot\vec b=|\vec a||\vec b|\cos\theta=a_xb_x+a_yb_y$（答えは数） & 4-12 \\
直交・なす角 & 直交なら $\vec a\cdot\vec b=0$。\ $\cos\theta=\dfrac{\vec a\cdot\vec b}{|\vec a||\vec b|}$ & 4-12 \\
正射影 & $\vec e$ の向きの成分 $=\vec e\cdot\vec b$（$\vec e$ は単位ベクトル） & 4-12 \\
空間の距離 & $\mathrm{AB}=\sqrt{(x_2-x_1)^2+(y_2-y_1)^2+(z_2-z_1)^2}$ & 4-13 \\
空間のベクトル & $\vec a=(a_x,\ a_y,\ a_z)=a_x\vec i+a_y\vec j+a_z\vec k$、\ $|\vec a|=\sqrt{a_x^2+a_y^2+a_z^2}$ & 4-13 \\
空間の内積 & $\vec a\cdot\vec b=a_xb_x+a_yb_y+a_zb_z$ & 4-13 \\
方向余弦・平面 & $\vec a\div|\vec a|=(\cos\theta_x,\ \cos\theta_y,\ \cos\theta_z)$。法線ベクトル $(a,\ b,\ c)$ の平面は $ax+by+cz=d$ & 4-13 \\
力のモーメント & $M=rF\sin\theta$。平面では $M=xF_y-yF_x$（左回りが＋） & 4-14 \\
外積 & 大きさ $|\vec a\times\vec b|=|\vec a||\vec b|\sin\theta$（平行四辺形の面積）。向きは $\vec a$ から $\vec b$ へ回す右ねじの進む向き & 4-14 \\
外積の成分 & $\vec a\times\vec b=(a_yb_z-a_zb_y,\ a_zb_x-a_xb_z,\ a_xb_y-a_yb_x)$。\ $\vec b\times\vec a=-\vec a\times\vec b$ & 4-14 \\
物理の外積 & トルク $\vec r\times\vec F$、\ 角運動量 $\vec r\times m\vec v$、\ 磁場からの力 $q\vec v\times\vec B$ & 4-14 \\
極座標 & $(r,\ \theta)$ の点は $x=r\cos\theta$、$y=r\sin\theta$。\ $r=\sqrt{x^2+y^2}$（$\theta$ は区画を見て決める） & {4-2\\4-8\\4-15} \\
単位ベクトル・円運動 & $\vec e_r=(\cos\theta,\ \sin\theta)$、\ $\vec e_\theta=(-\sin\theta,\ \cos\theta)$。速度 $\vec v=r\omega\,\vec e_\theta$ & 4-15 \\
円柱座標 & $(r,\ \theta,\ z)$：\ $x=r\cos\theta$、$y=r\sin\theta$、$z=z$ & 4-15 \\
球座標 & $(r,\ \theta,\ \varphi)$：\ $x=r\sin\theta\cos\varphi$、$y=r\sin\theta\sin\varphi$、$z=r\cos\theta$ & 4-15 \\
緯度と経度 & $\theta=\ang{90}-$（北緯）、\ $\varphi=$（東経） & 4-15 \\
\end{longtblr}\par}

\bksubsection{6}{電卓・表計算の使い方}
{\small
\begin{longtblr}{colspec={Q[l,m,wd=26mm] Q[l,m,wd=48mm] Q[l,m,wd=40mm] Q[c,m,wd=12mm]}, rowhead=1, row{1}={bg=bkVol!14, font=\sffamily\bfseries\footnotesize}, hlines={0.4pt,bkLine}, rowsep=1.6pt, colsep=4pt}
すること & 書き方・キー（例） & 意味 & 章 \\
角の単位 & 電卓の DEG・RAD & 度で計算するか、ラジアンで計算するか & {4-2\\4-6} \\
三角比から角 & $\sin^{-1}$・$\cos^{-1}$・$\tan^{-1}$（\texttt{SHIFT} を押してから） & $\tan^{-1}2$ は「$\tan$ が2になる角」 & 4-6 \\
円周率 & \texttt{=PI()} & $\pi=3.14159\cdots$ & 4-2 \\
度→ラジアン & \texttt{=RADIANS(A2)} & 180なら3.1416 & 4-2 \\
ラジアン→度 & \texttt{=DEGREES(B2)} & 3.1416なら180 & 4-2 \\
度のサイン & \texttt{=SIN(RADIANS(A2))} & \texttt{SIN} はラジアンで計算する & {4-2\\4-9} \\
点の向きの角 & \texttt{=ATAN2(}$x$\texttt{,} $y$\texttt{)} & 区画まで考えた角（ラジアン） & 4-15 \\
外積の成分 & \texttt{=C2*D3-D2*C3} など & 成分の式をセルで計算 & 4-14 \\
\end{longtblr}\par}
\begin{itemize}[leftmargin=1.4\zw, itemsep=2pt, topsep=4pt, label={\color{\bkTT}●}]
\item RAD のまま $\sin30$ を計算すると $-0.988$ になります。これは「30ラジアン」のサインです。答えがおかしいときは、まず DEG・RAD を確かめます（4-6）。
\item $\sin\theta=0.5$ になる角は、\ang{0}から\ang{180}のあいだに\ang{30}と\ang{150}の2つあります。電卓の $\sin^{-1}$ は1つしか出さないので、図をかいて確かめます（4-8）。
\item \texttt{ATAN2} は $x$ を先に入れます。プログラムの言語では $y$ を先に入れるものが多いので注意します（4-15）。
\end{itemize}

%% ---------------- チェック表 ----------------
\bkbacktitle{できるようになったことチェック表}
{\small 各章の「ふりかえり」を集めた表です。章を学び終えたら1回め、次の巻に進む前に2回めのチェックをしよう。できなかった所は、ページにもどって読みなおします。}\par\vspace{2mm}
%% 自動生成（tools/mkchecklist.py）。各章の「ふりかえり」の項目を集めたもの
\begingroup\footnotesize
\def\bkcb{\tikz[baseline=-0.5ex]\draw[\bkTF, line width=0.7pt, rounded corners=0.5pt] (0,-0.3em) rectangle (0.8em,0.5em);}
\begin{longtblr}{colspec={Q[c,m,wd=9mm] X[l,m] Q[c,m,wd=10mm] Q[c,m,wd=9mm] Q[c,m,wd=9mm]}, width=\linewidth, rowhead=1, colsep=3pt,
  row{1}={bg=bkVol!14, font=\sffamily\bfseries\scriptsize}, hlines={0.4pt,bkLine}, rowsep=0.7pt}
章 & できるようになったこと & ページ & 1回め & 2回め \\
{\sffamily\bfseries\color{\bkTT}4-1} & 平行線の同位角・錯角から、角を求められる。 & \pageref{bk:4-1:heikou} & \bkcb & \bkcb \\
 & 内角の和が180°のわけと、多角形の角の和がわかる。 & \pageref{bk:4-1:naikaku} & \bkcb & \bkcb \\
 & 定義・公理・定理・法則のちがいを説明でき、合同条件・逆・反例を使える。 & \pageref{bk:4-1:teigi} & \bkcb & \bkcb \\
{\sffamily\bfseries\color{\bkTT}4-2} & 弦・弧・おうぎ形がわかり、弧の長さとおうぎ形の面積を、中心角に比例させて求められる。 & \pageref{bk:4-2:hirei} & \bkcb & \bkcb \\
 & 角をラジアンで表し、度と換算して、$s=r\theta$ と $S=\tfrac{1}{2}r^2\theta$ を使える。 & \pageref{bk:4-2:rad} & \bkcb & \bkcb \\
 & 接線が半径に垂直なわけを説明でき、DEG・RAD と極座標 $(r,\ \theta)$ を使える。 & \pageref{bk:4-2:sessen} & \bkcb & \bkcb \\
{\sffamily\bfseries\color{\bkTT}4-3} & 拡大図・縮図と縮尺を使って、本当の長さや図の上の長さを求められる。 & \pageref{bk:4-3:kakudai} & \bkcb & \bkcb \\
 & 三角形の相似条件を見つけ、相似比や平行線と線分の比から、辺の長さを求められる。 & \pageref{bk:4-3:jouken} & \bkcb & \bkcb \\
 & 面積の比は $k^2$、体積の比は $k^3$ を使い、すいの断面や直角三角形の辺の比を説明できる。 & \pageref{bk:4-3:menseki} & \bkcb & \bkcb \\
{\sffamily\bfseries\color{\bkTT}4-4} & 方眼の面積から $a^2+b^2=c^2$ を見つけ、2通りに説明できる。 & \pageref{bk:4-4:shoumei} & \bkcb & \bkcb \\
 & 辺の長さ・2点の距離・直方体の対角線を求められる。 & \pageref{bk:4-4:kyori} & \bkcb & \bkcb \\
 & 直角三角形かを見分け、円の式のわけがわかる。 & \pageref{bk:4-4:en} & \bkcb & \bkcb \\
{\sffamily\bfseries\color{\bkTT}4-5} & 平行四辺形・三角形・台形・ひし形の面積を、長方形に直して求められる。 & \pageref{bk:4-5:sankaku} & \bkcb & \bkcb \\
 & 円の面積と、柱・すい・球の体積を、わけといっしょに使える。 & \pageref{bk:4-5:en} & \bkcb & \bkcb \\
 & 展開図から表面積を求め、体積を L に換算できる。 & \pageref{bk:4-5:tenkai} & \bkcb & \bkcb \\
{\sffamily\bfseries\color{\bkTT}4-6} & 直角三角形で、どの角から見るかに気をつけて、$\sin$・$\cos$・$\tan$ を求められる。 & \pageref{bk:4-6:teigi} & \bkcb & \bkcb \\
 & 三角比と電卓で、辺の長さや角度を求められる。坂の傾きを角度に直せる。 & \pageref{bk:4-6:atai} & \bkcb & \bkcb \\
 & 斜面の重力を、$mg\sin\theta$ と $mg\cos\theta$ に分けられる。 & \pageref{bk:4-6:bunkai} & \bkcb & \bkcb \\
{\sffamily\bfseries\color{\bkTT}4-7} & 鈍角の三角比を使い、三角形の面積を $\tfrac{1}{2}ab\sin C$ で求められる。 & \pageref{bk:4-7:menseki} & \bkcb & \bkcb \\
 & 正弦定理と余弦定理を使い分けて、三角形の辺や角を求められる。 & \pageref{bk:4-7:seigen} & \bkcb & \bkcb \\
 & ロボットのうでの関節の角を、余弦定理で求められる。 & \pageref{bk:4-7:robot} & \bkcb & \bkcb \\
{\sffamily\bfseries\color{\bkTT}4-8} & 単位円の上の点 $(\cos\theta,\ \sin\theta)$ を求められる。 & \pageref{bk:4-8:tanien} & \bkcb & \bkcb \\
 & $\sin^2\theta+\cos^2\theta=1$ や90°回す式を説明できる。 & \pageref{bk:4-8:sanheihou} & \bkcb & \bkcb \\
 & $\omega$ と $T$ で、回る点の位置・速さを求められる。 & \pageref{bk:4-8:omega} & \bkcb & \bkcb \\
{\sffamily\bfseries\color{\bkTT}4-9} & 加法定理を単位円で導き、$\sin(a\pm b)$・$\cos(a\pm b)$ を計算できる。 & \pageref{bk:4-9:kahou} & \bkcb & \bkcb \\
 & 2倍角の式と $\sin^2\theta=\tfrac{1-\cos2\theta}{2}$ を作り、使える。 & \pageref{bk:4-9:nibai} & \bkcb & \bkcb \\
 & 合成と「積と和の式」で、うなりや電力を説明できる。 & \pageref{bk:4-9:gousei} & \bkcb & \bkcb \\
{\sffamily\bfseries\color{\bkTT}4-10} & 向きと大きさをもつ量を矢印と記号で表せる。 & \pageref{bk:4-10:kigou} & \bkcb & \bkcb \\
 & 矢印をつないで、合力の大きさを求められる。 & \pageref{bk:4-10:wa} & \bkcb & \bkcb \\
 & 差をかいて、速度の変化や相対速度を求められる。 & \pageref{bk:4-10:sa} & \bkcb & \bkcb \\
{\sffamily\bfseries\color{\bkTT}4-11} & 成分で、和・差・実数倍・大きさを計算できる。 & \pageref{bk:4-11:keisan} & \bkcb & \bkcb \\
 & 力や速度を、成分に分けて考えられる。 & \pageref{bk:4-11:bunkai} & \bkcb & \bkcb \\
 & 内分点・重心・進む点の位置を求められる。 & \pageref{bk:4-11:ichi} & \bkcb & \bkcb \\
{\sffamily\bfseries\color{\bkTT}4-12} & 仕事を求め、$W=Fs\cos\theta$ を使える。 & \pageref{bk:4-12:naname} & \bkcb & \bkcb \\
 & 内積を、大きさと角からも、成分からも計算できる。 & \pageref{bk:4-12:naiseki} & \bkcb & \bkcb \\
 & 内積で、なす角・直交・相関係数を求められる。 & \pageref{bk:4-12:chokkou} & \bkcb & \bkcb \\
{\sffamily\bfseries\color{\bkTT}4-13} & 空間の点を $(x,\ y,\ z)$ で表し、2点の距離を求められる。 & \pageref{bk:4-13:zahyo} & \bkcb & \bkcb \\
 & 3次元のベクトルを成分で計算し、内積からなす角を求められる。 & \pageref{bk:4-13:naiseki} & \bkcb & \bkcb \\
 & 基本ベクトル・単位ベクトル・法線ベクトルで、向きや平面を表せる。 & \pageref{bk:4-13:kihon} & \bkcb & \bkcb \\
{\sffamily\bfseries\color{\bkTT}4-14} & 力のモーメントを「力×うでの長さ」で求め、$M=xF_y-yF_x$ で計算できる。 & \pageref{bk:4-14:moment} & \bkcb & \bkcb \\
 & 外積の向きと大きさがわかり、成分で計算できる。 & \pageref{bk:4-14:gaiseki} & \bkcb & \bkcb \\
 & トルクや、磁場から受ける力の向きを、外積で求められる。 & \pageref{bk:4-14:torque} & \bkcb & \bkcb \\
{\sffamily\bfseries\color{\bkTT}4-15} & 極座標と $(x,\ y)$ を、区画に注意して行き来できる。 & \pageref{bk:4-15:kyoku} & \bkcb & \bkcb \\
 & $\vec e_r$、$\vec e_\theta$ を使って、円運動の速度を $r\omega\,\vec e_\theta$ と表せる。 & \pageref{bk:4-15:er} & \bkcb & \bkcb \\
 & 円柱座標・球座標で空間の点を表し、$(x,\ y,\ z)$ に直せる。 & \pageref{bk:4-15:enchu} & \bkcb & \bkcb \\
\end{longtblr}
\endgroup


%% ---------------- つぎの巻へ（第5巻の予告と、この巻の道具がこの先どこで使われるか） ----------------
\bkbacktitle{つぎの巻へ}
\newcommand\bknextpill[1]{\bk@pill{#1}}
\def\bknextrow{1}
\newcommand\bknextcard[4]{%
  \begin{tcolorbox}[enhanced, nobeforeafter, box align=top, width=0.485\linewidth, colframe=bkVol5!75!black, colback=white, boxrule=0.9pt, arc=2.5mm,
    left=2.5mm, right=2.5mm, top=1.5mm, bottom=2mm, equal height group=bknext\bknextrow]
  \bknextpill{#1}\enspace{\sffamily\bfseries\small #2}\par\vspace{1mm}
  {\sffamily\bfseries\footnotesize\color{bkDeep5}#3}\par\vspace{1mm}
  {\footnotesize\setlength{\parindent}{0pt}{\sffamily\bfseries\color{bkInk}使う道具}\enspace #4}
  \end{tcolorbox}}
長さと角度で、形と向きを数にする道具がそろいました。第5巻「変化をつかむ」では、この巻の図形・三角関数・ベクトルを使って、変わる速さ（微分）と、たまった量（積分）を計算します。第5巻の章をのぞいてみましょう。\par\vspace{2.5mm}
\noindent\bknextcard{5-3}{かぎりなく近づく：極限と接線}{曲線にぴったりふれる直線の傾き}{円の接線（4-2）と、ラジアンで小さな角を計算すること（4-2）。}\hfill
\bknextcard{5-9}{グラフの面積は進んだ距離}{速さのグラフの面積が、進んだ道のりになる}{細かく分けてたす考え方と、台形の面積（4-5）。}\par\vspace{1.2mm}
\def\bknextrow{2}\noindent\bknextcard{5-10}{細い長方形に分けてたそう}{曲がった形の面積を、細い長方形で求める}{面積と体積を細かく分けてたす（4-5）と、相似（4-3）。}\hfill
\bknextcard{5-12}{物理での積分：仕事と流量}{ばねをのばす仕事は、三角形の面積}{仕事＝力×力の向きに動いた距離と、内積（4-12）。}\par\vspace{1.2mm}
\def\bknextrow{3}\noindent\bknextcard{5-15}{sin・cosを微分すると？}{微分すると、波が1/4周期ずれる}{三角関数と円運動（4-8）、ラジアン（4-2）、加法定理（4-9）。}\hfill
\bknextcard{5-17}{近くでは多項式で近似できる：テイラー展開}{小さくゆれる振り子は、$\sin\theta\approx\theta$で計算できる}{ラジアンの角（4-2）と、三角比（4-6）。}\par
\bkfukidashi{delta}{第4巻の図形・三角関数・ベクトルは、この先の全部の巻で使うよ。チェック表で自信のない章を見なおしてから、第5巻へ進もう！}

\clearpage
\bkbkhead{この巻の道具は、この先でこう使う}
{\small 第4巻の章と、その道具を使う第5巻から先の章です（章のはじめの「道しるべ」の「この章の先に」から集めました）。}\par\vspace{3mm}
%% 自動生成（tools/mknext.py）。各章の道しるべ「この章の先に」から、後ろの巻の章を集めたもの
\begin{longtblr}{colspec={Q[l,m,wd=58mm] X[l,m]}, width=\linewidth, rowhead=1, row{1}={bg=bkVol!14, font=\sffamily\bfseries\footnotesize},
  hlines={0.4pt,bkLine}, rowsep=1.5pt, cells={font=\footnotesize}}
この巻の章 & この道具を使う章 \\
{\bknextpill{4-1}\ きまりから考えを積み上げる：定義・公理・定理} & {\bknextpill{6-15}\ まとめ：物理は「微分方程式＋はじめの状態」で書ける} \\
{\bknextpill{4-2}\ 角度と円、ラジアン} & {\bknextpill{5-3}\ かぎりなく近づく：極限と接線\\ \bknextpill{5-15}\ sin・cosを微分すると？\\ \bknextpill{5-17}\ 近くでは多項式で近似できる：テイラー展開} \\
{\bknextpill{4-3}\ にた形と比：相似} & {\bknextpill{5-10}\ 細い長方形に分けてたそう} \\
{\bknextpill{4-4}\ 三平方の定理} & {\bknextpill{5-16}\ 積分の道具箱：置換積分と部分積分} \\
{\bknextpill{4-5}\ 面積と体積：細かく分けてたす} & {\bknextpill{5-4}\ 表から導関数を見つけよう\\ \bknextpill{5-6}\ 微分の道具箱：積と合成関数\\ \bknextpill{5-9}\ グラフの面積は進んだ距離\\ \bknextpill{5-10}\ 細い長方形に分けてたそう} \\
{\bknextpill{4-6}\ 坂道と三角比} & {\bknextpill{5-17}\ 近くでは多項式で近似できる：テイラー展開} \\
{\bknextpill{4-8}\ 三角関数：円運動の影} & {\bknextpill{5-15}\ sin・cosを微分すると？\\ \bknextpill{5-16}\ 積分の道具箱：置換積分と部分積分\\ \bknextpill{6-6}\ ばねと振り子\\ \bknextpill{7-6}\ 回転行列：くるっと回す} \\
{\bknextpill{4-9}\ 波をたし合わせる：加法定理とうなり} & {\bknextpill{5-15}\ sin・cosを微分すると？\\ \bknextpill{6-10}\ オイラーの公式\\ \bknextpill{9-1}\ 波をたし合わせて形を作る：フーリエ級数} \\
{\bknextpill{4-10}\ 向きと大きさをもつ量：ベクトル} & {\bknextpill{7-2}\ 行列×ベクトル：形を変える機械\\ \bknextpill{8-5}\ ベクトル場} \\
{\bknextpill{4-11}\ ベクトルを成分で計算しよう} & {\bknextpill{5-15}\ sin・cosを微分すると？\\ \bknextpill{7-2}\ 行列×ベクトル：形を変える機械} \\
{\bknextpill{4-12}\ 内積：どれだけ同じ向きか} & {\bknextpill{5-12}\ 物理での積分：仕事と流量\\ \bknextpill{8-7}\ 道にそってたす：線積分\\ \bknextpill{9-2}\ 成分の取り出し方：直交性} \\
{\bknextpill{4-13}\ 空間の座標と3次元のベクトル} & {\bknextpill{7-7}\ 3次元の回転と座標の変換\\ \bknextpill{8-5}\ ベクトル場} \\
{\bknextpill{4-14}\ 外積：回す力と右ねじ} & {\bknextpill{7-4}\ 行列式：面積が何倍になるか\\ \bknextpill{8-10}\ うず：回転 rot} \\
{\bknextpill{4-15}\ いろいろな座標：極座標・円柱座標・球座標} & {\bknextpill{8-6}\ 面や立体でたし合わせる：重積分\\ \bknextpill{8-12}\ ならす働き：ラプラシアン} \\
\end{longtblr}


%% ---------------- さくいん ----------------
\bkbacktitle{さくいん}
{\small 数字はページです。その語がはじめて出てくるページや、くわしく説明しているページをのせています。}\par\vspace{2mm}
\begin{multicols}{2}\small\setlength{\parskip}{0pt}
\directlua{bk.printidx()}
\end{multicols}

%% ---------------- 奥付 ----------------
\bkpadfour
\global\bk@decofalse\bkfolio{false}
\null\vfill
\noindent\begin{minipage}{\linewidth}\small
{\sffamily\bfseries\Large\color{bkDeep}物理を理解するための数学　第4巻}\par\vspace{1mm}
{\sffamily\bfseries\large 形と向き}\par\vspace{3mm}
{\color{bkVol}\hrule height 1pt}\vspace{3mm}
\begin{tabular}{@{}l@{\hspace{6mm}}l@{}}
編者 & 「物理を理解するための数学」編集委員会 \\
 & 委員長　小学校長（まとめ役） \\
 & 委員　数学・電気工学・機械工学・教育学・量子力学・建築工学・\\
 & \phantom{委員　}半導体工学の大学教授、幼稚園教諭、塾講師、\\
 & \phantom{委員　}特別支援学校教諭、心理士、科学解説者、小学校教諭 \\
デザイン & ユニバーサルデザイン担当委員 \\
挿絵 & イラスト担当委員 \\
版 & 第1版（第13回編集委員会）　2026年10月 \\
\end{tabular}\par\vspace{3mm}
{\color{bkVol}\hrule height 1pt}\vspace{2mm}
{\footnotesize この本の問題の答えは、巻末の「こたえ」にあります。くわしい解き方は、別冊の問題集と指導書にのせます。}
\end{minipage}
\makeatother
